
\subsubsection{3.1 Error Class Independence Measurement}

For a code generation benchmark with problems $P = \{p_1, \ldots, p_n\}$, each verification strategy $s$ produces an improvement set $I_s \subseteq P$ containing problems where strategy $s$ increases correctness.

Independence between strategies $s_1$ and $s_2$ is measured using the Jaccard index:

$$J(s_1, s_2) = \frac{|I_{s_1} \cap I_{s_2}|}{|I_{s_1} \cup I_{s_2}|}$$

Interpretation:
\begin{itemize}
\item $J = 0$: Complete independence (disjoint improvement sets)
\item $J = 1$: Complete overlap (identical improvement sets)
\item $J < 0.30$: Largely independent (threshold for multiplicative model)
\end{itemize}

\subsubsection{3.2 Strategy Definitions}

Three verification strategies are defined corresponding to different error classes:

\textbf{Grammar Constraints ($s_g$):} Token-level logit masking via deterministic finite automaton (DFA). Uses SynCode in grammar\_strict mode with Python grammar. A problem is counted as improved if it passes after constrained generation when it failed with unconstrained baseline.

\textbf{Static Analysis ($s_a$):} AST-level pattern matching using Bandit (security) and Pylint (reliability). Improvement is measured as reduction in static analysis issues after feedback-based refinement.

\textbf{SMT Verification ($s_m$):} Formal specification checking using Z3 on HumanEval-Verus problems (23 problems with Verus specifications). Improvement requires code satisfying formal specification after SMT-guided repair.

\subsubsection{3.3 Pipeline Architecture}

The designed pipeline is ordered by abstraction level:

\begin{verbatim}
[Raw LLM Code]
    ↓
[Stage 1: Grammar Constraints]
    - Mechanism: Token-level logit masking via DFA
    - Targets: Syntax errors (compilation failures)
    ↓
[Stage 2: Static Analysis]
    - Mechanism: AST-based pattern matching + LLM repair
    - Targets: Semantic issues (security, reliability)
    ↓
[Stage 3: SMT Verification]
    - Mechanism: Z3 constraint solving + guided repair
    - Targets: Specification violations
\end{verbatim}

\subsubsection{3.4 Hypothesis Structure}

The main hypothesis (multiplicative error reduction) is decomposed into testable sub-hypotheses:

\begin{itemize}
\item \textbf{H-E1 (Existence):} Error classes are largely independent (Jaccard < 0.30 for all strategy pairs).
\item \textbf{H-M1 (Mechanism):} Grammar-constrained decoding reduces syntax errors through prefix automata enforcement.
\item \textbf{H-M2 (Mechanism):} Static analysis feedback reduces semantic issues when using instruction-tuned models.
\item \textbf{H-M3 (Mechanism):} SMT-guided repair enforces specification satisfaction on semantically-filtered code.
\item \textbf{H-M4 (Mechanism):} Pipeline composition achieves synergy coefficient $S \geq 0.8$, where $S = \text{Combined improvement} / \sum \text{Individual improvements}$.
\end{itemize}
